\begin{lemma}
\label{editorial-source-lemma-finite-finite-fibres}
Suppose $R \to S$ is finite.
Then the fibres of $\Spec(S)\to\Spec(R)$ are finite.
More precisely, if $s_1,\ldots,s_m$ generate $S$ as an $R$-module,
then for every $\mathfrak p\in\operatorname{Spec}(R)$,
$$
\begin{aligned}
\#\{\mathfrak q:\varphi^{-1}(\mathfrak q)=\mathfrak p\}
&\leq\sum_{\varphi^{-1}(\mathfrak q)=\mathfrak p}
[\kappa(\mathfrak q):\kappa(\mathfrak p)]
\\
&\leq\dim_{\kappa(\mathfrak p)}(S\otimes_R\kappa(\mathfrak p))
\leq m.
\end{aligned}
$$
The middle difference is exactly the dimension of the nilradical
of this fibre algebra. In particular the residue-degree sum equals
the fibre dimension exactly when this fibre is reduced.
\end{lemma}

\begin{proof}
By the discussion in
Remark \ref{algebra-remark-fundamental-diagram}
the fibres are the spectra of the rings $S \otimes_R \kappa(\mathfrak p)$.
As $R \to S$ is finite, these fibre rings are finite over
$\kappa(\mathfrak p)$ hence Noetherian by
Lemma \ref{editorial-source-lemma-Noetherian-permanence}.
By
Lemma \ref{editorial-source-lemma-integral-no-inclusion}
every prime of $S \otimes_R \kappa(\mathfrak p)$ is a minimal
prime. Hence by
Lemma \ref{editorial-source-lemma-Noetherian-irreducible-components}
there are at most finitely many.
\medskip\noindent
For the quantitative assertions, keep the original map
$\varphi:R\to S$ and the given module generators $s_1,\ldots,s_m$.
Fix $\mathfrak p\in\operatorname{Spec}(R)$ and put
$K=\kappa(\mathfrak p)$, $U=R\setminus\mathfrak p$,
$A=S\otimes_R K$ and $d=\dim_K A$.
The original surjection $R^m\to S$ induces the surjection
$$
K^m\longrightarrow A,\qquad
(\lambda_j)_j\longmapsto\sum_{j=1}^m s_j\otimes\lambda_j.
$$
Indeed $s=\sum_j\varphi(a_j)s_j$ gives
$s\otimes\lambda=\sum_j s_j\otimes(\overline{a_j}\lambda)$,
and every tensor is a finite sum of pure tensors. Thus $d\leq m$.

The exact fibre-ring isomorphism and its inverse are
$$
\begin{aligned}
A&\longrightarrow U^{-1}S/\mathfrak pU^{-1}S,&
s\otimes\frac{r+\mathfrak p}{u+\mathfrak p}
&\longmapsto\frac{s\varphi(r)}{\varphi(u)}+\mathfrak pU^{-1}S,\\
\frac{s}{\varphi(u)}+\mathfrak pU^{-1}S
&\longmapsto s\otimes\frac1{u+\mathfrak p}.&&
\end{aligned}
$$
Here $U^{-1}S$ means localization at the image $\varphi(U)$.
Tensor balancing respects the original coefficient map; a change
of fraction representatives is killed by a denominator from $U$
on both sides, and an element of $\mathfrak p$ acts as zero on
$K$. These facts verify the defining relations of both maps.
Their composites on the displayed fractions and pure tensors
are identities, proving the isomorphism.
Write $\beta:S\to A$, $s\mapsto s\otimes1$.
A prime $\mathfrak q$ of $S$ with
$\varphi^{-1}(\mathfrak q)=\mathfrak p$ corresponds to
$$
\mathfrak n_{\mathfrak q}
=U^{-1}\mathfrak q/\mathfrak pU^{-1}S\subset A.
$$
It is proper because $\mathfrak q$ avoids $\varphi(U)$, and
contraction recovers $\mathfrak q$: the condition
$s/1\in U^{-1}\mathfrak q$ is equivalent to
$\varphi(u)s\in\mathfrak q$ for some $u\in U$, hence to
$s\in\mathfrak q$ by primality. Conversely, a prime
$\mathfrak n\subset A$ contracts to a prime $\mathfrak q\subset S$
whose contraction in $R$ is $\mathfrak p$: elements of
$\mathfrak p$ act as zero and elements of $U$ act as units.
Every element of $A$ is represented by a fraction $s/\varphi(u)$
modulo $\mathfrak p$, and it belongs to $\mathfrak n$ precisely
when $\beta(s)$ does. This verifies the remaining composite of
the two prime correspondences. In particular the correspondence
is a bijection for the original, possibly noninjective map.

For every prime $\mathfrak n$ of $A$, the quotient $A/\mathfrak n$
is a finite dimensional domain over $K$, hence a field by
Lemma \ref{editorial-source-lemma-integral-over-field}. Therefore all these primes
are maximal. If $\mathfrak n=\mathfrak n_{\mathfrak q}$,
its field is canonically the original residue field:
$$
A/\mathfrak n_{\mathfrak q}
\cong U^{-1}(S/\mathfrak q)
\xrightarrow{\ \sim\ }\kappa(\mathfrak q),\qquad
\frac{s+\mathfrak q}{\varphi(u)+\mathfrak q}
\longmapsto\frac{s+\mathfrak q}{\varphi(u)+\mathfrak q}.
$$
The last map is injective since $S/\mathfrak q$ is a domain.
Its source is already a field containing $S/\mathfrak q$,
so it contains the inverse of each nonzero element of that domain;
it therefore contains every fraction and the map is surjective.
This also identifies the coefficient field $K$ in both targets.

Here is the Chinese remainder map with its actual preimages.
For any finite family of distinct maximal ideals
$\mathfrak n_1,\ldots,\mathfrak n_t$ of $A$, comaximality gives
$b_{ij}\in\mathfrak n_j$ with $1-b_{ij}\in\mathfrak n_i$
when $i\ne j$. Put $e_i=\prod_{j\ne i}b_{ij}$.
It is $1$ modulo $\mathfrak n_i$ and $0$ modulo each other
$\mathfrak n_j$. Hence
$$
\rho:A\longrightarrow\prod_{i=1}^t A/\mathfrak n_i,
\qquad a\longmapsto(a+\mathfrak n_i)_i
$$
is surjective: representatives $a_i$ of a prescribed tuple have
preimage $\sum_i e_ia_i$. Its kernel is exactly
$\bigcap_i\mathfrak n_i$. Comparing $K$-dimensions gives
$\sum_i\dim_K(A/\mathfrak n_i)\leq d$. Every field factor has
positive dimension, so there cannot be more than $d$ maximal
ideals: any $d+1$ distinct ones would contradict the inequality.
If $A\ne0$ it has a maximal ideal, for example a proper ideal of
greatest $K$-dimension. Thus its primes are a finite family and
the preceding construction applies to all of them at once.
Let $L_i=A/\mathfrak n_i\cong\kappa(\mathfrak q_i)$ and
$f_i=[L_i:K]$. With $J=\bigcap_i\mathfrak n_i$, it gives
$$
0\longrightarrow J\longrightarrow A\xrightarrow{\rho}
\prod_iL_i\longrightarrow0,
\qquad
d=\sum_i f_i+\dim_KJ.
$$

The ideal $J$ is precisely the nilradical, and is nilpotent.
To prove the latter assertion directly, the decreasing chain
$J,J^2,J^3,\ldots$ stabilizes since these are subspaces of the
finite dimensional $K$-space $A$. Choose $N\geq1$ with
$J^N=J^{N+1}$. A finite $K$-basis of $J^N$ generates it as an
$A$-module. Write each generator as a $J$-linear combination
of all generators, giving a matrix $C$ with entries in $J$.
Multiplication by the adjugate of $I-C$ gives
$\det(I-C)J^N=0$. This determinant belongs to $1+J$ and is a
unit: if it lay in a maximal ideal, that ideal would also contain
$J$, and hence would contain $1$. Therefore $J^N=0$.
Every nilpotent element lies in all prime ideals, and conversely
every element of $J$ is now nilpotent, proving the assertion.
Consequently the middle inequality is an equality exactly when
$A$ is reduced. Together with $f_i\geq1$ and $d\leq m$, this
proves all the displayed inequalities. For $A=0$ there are no
prime ideals, its dimension is zero, and the formulas use empty sums.

There is a more precise local formula, retaining every multiplicity.
For $A\ne0$ and the above $N$, the ideals $\mathfrak n_i^N$ are
pairwise comaximal. In fact, if $a+b=1$ with
$a\in\mathfrak n_i$, $b\in\mathfrak n_j$, every term of
$(a+b)^{2N-1}=1$ belongs to $\mathfrak n_i^N$ or $\mathfrak n_j^N$.
For two comaximal ideals their intersection is their product:
if $x$ is in both and $a+b=1$ with one element in each ideal,
then $x=xa+xb$ belongs to the product. Induction gives the finite
version; the product of the first ideals stays comaximal with
the next since multiplying equations congruent to $1$ modulo
the next ideal gives another such equation. It follows that
$$
\bigcap_i\mathfrak n_i^N
=\prod_i\mathfrak n_i^N
=\left(\prod_i\mathfrak n_i\right)^N
=J^N=0.
$$
The same Chinese remainder construction therefore gives
$$
A\xrightarrow{\ \sim\ }\prod_i C_i,
\qquad C_i=A/\mathfrak n_i^N.
$$
Each $C_i$ has unique maximal ideal
$\mathfrak m_i=\mathfrak n_i/\mathfrak n_i^N$: a prime
containing $\mathfrak n_i^N$ contains every element of
$\mathfrak n_i$, by taking its $N$th power, and thus equals
that maximal ideal. The residue field is $L_i$.

The identification $C_i\cong A_{\mathfrak n_i}$ has explicit maps.
Choose the Chinese remainder element $e_i$ that is $1$ modulo
$\mathfrak n_i^N$ and $0$ modulo the other powers. For
$a\in\mathfrak n_i^N$, the product $e_ia$ is in the intersection
of all powers and is zero. Since $e_i\notin\mathfrak n_i$,
this gives $a/1=0$ in $A_{\mathfrak n_i}$. Thus
$\overline a\mapsto a/1$ defines $C_i\to A_{\mathfrak n_i}$.
Its inverse sends $a/u$ to $\overline a\,\overline u^{-1}$;
the inverse exists since $u\notin\mathfrak n_i$ and $C_i$ is
local. The two composites on elements and fractions are identities.
There is also the exact original-ring comparison
$$
A_{\mathfrak n_i}\cong S_{\mathfrak q_i}/\mathfrak pS_{\mathfrak q_i}.
$$
One direction sends $s/v$ modulo $\mathfrak p$ to
$\beta(s)/\beta(v)$ for $v\notin\mathfrak q_i$.
The other is induced by
$s\otimes((r+\mathfrak p)/(u+\mathfrak p))
\mapsto s\varphi(r)/\varphi(u)$ modulo $\mathfrak p$.
It extends to localization at $\mathfrak n_i$, since an element
outside that maximal ideal maps to an element with nonzero
residue in the identified field $\kappa(\mathfrak q_i)$ and
hence to a unit of the target local ring. Tensor, fraction and
quotient relations are respected as in the preceding fibre map,
and the composites fix every such fraction, proving the comparison.

Use the finite filtration
$C_i\supset\mathfrak m_i\supset\mathfrak m_i^2\supset\cdots\supset0$.
Each quotient $\mathfrak m_i^j/\mathfrak m_i^{j+1}$ is annihilated
by $\mathfrak m_i$ and is a finite dimensional $L_i$-space.
Put
$$
\mu_i=\sum_{j\geq0}
\dim_{L_i}(\mathfrak m_i^j/\mathfrak m_i^{j+1}).
$$
This positive integer is the length of $C_i$ over itself:
refine each quotient by a vector-space flag, and note that every
simple module over the local ring $C_i$ is $L_i$, because the
annihilator of a nonzero simple-module generator is its unique
maximal ideal. Every composition series has the same number of
factors, since each such factor has $K$-dimension $f_i$.
Taking dimensions of the filtration and of the product yields
the exact formulas
$$
\begin{aligned}
\dim_K(S\otimes_RK)&=\sum_i\mu_i[\kappa(\mathfrak q_i):K],\\
\dim_KJ&=\sum_i(\mu_i-1)[\kappa(\mathfrak q_i):K],\\
\mu_i&=\operatorname{length}_{S_{\mathfrak q_i}/\mathfrak pS_{\mathfrak q_i}}
 (S_{\mathfrak q_i}/\mathfrak pS_{\mathfrak q_i}).
\end{aligned}
$$
Thus no local multiplicity has been absorbed into a residue degree.

The empty fibre also has an exact criterion:
$A=0$ if and only if $\ker\varphi\not\subseteq\mathfrak p$.
If some $u\in U$ is in the kernel, its zero image is inverted
in $U^{-1}S$, giving the zero ring. Conversely $A=0$ implies
$U^{-1}S=\mathfrak pR_{\mathfrak p}(U^{-1}S)$.
This is a finite $R_{\mathfrak p}$-module generated by the images
of the specified $s_j$. Nakayama's Lemma \ref{algebra-lemma-NAK} gives
$U^{-1}S=0$, so $1/1=0$ implies $\varphi(u)=0$ for some
$u\in U$. This proves the criterion without assuming injectivity.

For clarity, all equality cases follow from these formulas.
The number of primes equals the residue-degree sum precisely
when every $f_i=1$. That sum equals $d$ precisely when $J=0$,
equivalently all $\mu_i=1$. The upper equality $d=m$ holds
precisely when the original images $s_j\otimes1$ are a $K$-basis.
In particular the number of primes equals $d$ exactly when
$A\cong K^t$, where $t$ is the number of primes: the equalities
give $J=0$, $f_i=1$ and the proved Chinese remainder isomorphism;
the converse follows from that product's coordinate maximal ideals.
The number of primes equals $m$ exactly when $A\cong K^m$.
These statements include the empty product, interpreted as the zero ring.
In terms of the original relation module
$H=\ker(R^m\to S)$, the kernel of $K^m\to A$ is exactly
$$
\frac{H_{\mathfrak p}+\mathfrak pR_{\mathfrak p}^m}
{\mathfrak pR_{\mathfrak p}^m},
$$
because localization gives $S_{\mathfrak p}=R_{\mathfrak p}^m/H_{\mathfrak p}$
and reduction gives that quotient. Thus $d=m$ exactly when
$H_{\mathfrak p}\subseteq\mathfrak pR_{\mathfrak p}^m$, equivalently
$H\subseteq\mathfrak pR^m$. The last equivalence is coordinatewise:
$a/1\in\mathfrak pR_{\mathfrak p}$ if and only if $a\in\mathfrak p$,
by the fraction equality criterion and primality.

The bounds and their limitations can be seen in explicit examples.
For $R=k$, $S=k^m$ and the coordinate idempotent generators,
the fibre has $m$ primes and all inequalities are equalities.
For $R=k[t]$, $S=k[x]$ and $\varphi(t)=x^m$ with $m\geq1$,
the original generators $1,x,\ldots,x^{m-1}$ form a free
$R$-basis: collecting powers by their residue modulo $m$ gives
spanning and uniqueness of the coefficients in $k[t]$.
At $\mathfrak p=(t)$ the fibre is $k[x]/(x^m)$, so it has
one prime, residue degree $1$, multiplicity $m$ and nilradical
dimension $m-1$, although both original rings are domains.
For the noninjective map $k[t]\to k$, $t\mapsto0$, with generator
$1$, the fibre at $(t-1)$ is zero, while the fibre at $(t)$ is $k$.
Repeating the generator $1$ for $R=S=k$ gives $d=1<m$ whenever
the supplied family has $m>1$ entries.

Equality in the residue-degree bound does not imply separability.
In characteristic $\ell>0$, let $K=\mathbb F_\ell(t)$ and
$L=\mathbb F_\ell(u)$ with $t\mapsto u^\ell$.
The elements $1,u,\ldots,u^{\ell-1}$ are linearly independent
over $K$: clearing denominators in any relation gives
$\sum_{j=0}^{\ell-1}u^jP_j(u^\ell)=0$, whose powers in distinct
residue classes modulo $\ell$ cannot cancel. Their span is closed
under multiplication since $u^\ell=t$, and is a finite dimensional
domain over $K$, hence a field. It contains $u$ and therefore is
the whole $L$, proving $[L:K]=\ell$ and
$L\cong K[X]/(X^\ell-t)$. Its unique residue field has degree
equal to its dimension and it is reduced. But
$$
L\otimes_KL\cong L[X]/(X^\ell-u^\ell)
=L[X]/((X-u)^\ell)
$$
is not reduced: $X-u$ is a nonzero element of nilpotence order
$\ell$, as its degree is smaller than that of the defining monic
polynomial. Thus the equality asserts reducedness of the specified
fibre, not reducedness after every field extension.

Finally every finite choice of positive integers $f_i,\mu_i$
is realized, showing the numerical formula is sharp. Take
$K=F(t_0)$ for a field $F$ and
$L_i=F(u_i)$ with $t_0\mapsto u_i^{f_i}$.
The same distinct-residue-class argument proves
$[L_i:K]=f_i$, with basis $1,u_i,\ldots,u_i^{f_i-1}$;
the spanning ring is a finite dimensional domain and hence a field
containing $u_i$, so the argument applies to the full rational
function field. Set
$$
S=\prod_i L_i[\varepsilon_i]/(\varepsilon_i^{\mu_i}).
$$
Its $K$-basis consists of the elements supported in factor $i$
with entries $u_i^a\varepsilon_i^b$ for
$0\leq a<f_i$, $0\leq b<\mu_i$. The nilpotent ideal in each
factor is $(\varepsilon_i)$, its quotient is the field $L_i$,
and its successive powers have one dimensional quotients over
$L_i$ until exponent $\mu_i$. Thus its unique prime has residue
degree $f_i$ and local length $\mu_i$. A prime of the product
selects exactly one factor because the coordinate idempotents
are orthogonal and sum to $1$. Consequently the fibre over the
zero prime of $K$ has exactly these degrees and multiplicities.
The displayed basis has size $\sum_i\mu_if_i$; appending zeros
gives a specified generating family of any larger size $m$.
Thus positivity and $\sum_i\mu_if_i\leq m$ are the complete
universal numerical restrictions on these data.
\end{proof}